% Generated draft: requires mathematical and source review.
\documentclass[11pt]{article}
\usepackage[margin=25mm]{geometry}
\usepackage{fontspec}
\setmainfont{Latin Modern Roman}
\usepackage{amsmath,amssymb,mathtools,unicode-math}
\setmathfont{Latin Modern Math}
\usepackage{enumitem,needspace,xurl,hyperref}
\setlength{\parindent}{0pt}
\setlength{\parskip}{5pt}
\setlength{\emergencystretch}{3em}
\newcommand{\N}{\mathbb{N}}
\newcommand{\Z}{\mathbb{Z}}
\newcommand{\Q}{\mathbb{Q}}
\newcommand{\R}{\mathbb{R}}
\newcommand{\C}{\mathbb{C}}
\newcommand{\F}{\mathbb{F}}
\newcommand{\A}{\mathbb{A}}
\newcommand{\PP}{\mathbb{P}}
\newcommand{\E}{\mathbb{E}}
\newcommand{\eps}{\varepsilon}
\newcommand{\dd}{\,\mathrm{d}}
\begin{document}
\section*{UnsolvedMath ID: 30001276}
\textbf{Source problem code:} catalogue entry
\subsection{Definitions and mathematical statement}
Let \(k\) be an imaginary quadratic field. The supplied formulation takes \(K/k\) to be a finite extension contained in a cyclotomic extension of \(k\), and takes \(F\) to be the maximal totally real subfield of \(K\). The intended cyclotomic extension and any additional conditions on \(K\) are not specified.

Let \(E\) denote the elliptic curve appearing in the source record, and let \(p\geq 5\) be a prime. Write \(E[p]\) for its group of geometric \(p\)-torsion points, a two-dimensional \(\F_p\)-vector space with its Galois action. The record does not specify the field of definition of \(E\) or the Galois group whose representation is intended. Let \(N\), \(\phi_F\), and the minimal modular degree denote the conductor or level, automorphic form, and modular degree intended by the source; their precise definitions and required modularity hypotheses are absent from the supplied extraction.

The stated hypotheses are:
\begin{enumerate}
\item the root number of \(L(E/k,s)\) is \(+1\);
\item the representation on \(E[p]\) has image \(\operatorname{GL}_2(\F_p)\);
\item \(p\) does not divide the minimal modular degree;
\item for every finite prime \(v\) of \(F\) satisfying
\[
v^2\mid N\mathcal O_F
\qquad\text{and}\qquad
p\mid\bigl(\operatorname{N}(v)+1\bigr),
\]
the restriction of \(E[p]\) to \(I_v\) is irreducible. Here \(\operatorname{N}(v)=\#(\mathcal O_F/v)\), and \(I_v\) denotes an inertia group at \(v\); the restriction presupposes a compatible Galois representation over \(F\), which the extraction does not specify.
\end{enumerate}

Write \(K[p^\infty]\) for the source's dihedral ring-class tower and \(J_K(\infty)\) for the group used in its Iwasawa algebra. Neither the tower nor this group is defined in the supplied record. For a profinite group \(H\), the conventional \(\Z_p\)-coefficient Iwasawa algebra is
\[
\Z_p[[H]]:=\varprojlim_U\Z_p[H/U],
\]
where \(U\) ranges over the open normal subgroups of \(H\). The notation \(\Lambda(J_K(\infty))\) is retained from the record, whose coefficient ring and treatment of finite torsion factors are unspecified.

The notation \(X(E/L)\) conventionally denotes the Pontryagin dual of the relevant \(p\)-primary Selmer group:
\[
X(E/L)=\operatorname{Hom}_{\mathrm{cont}}\!\left(\operatorname{Sel}_{p^\infty}(E/L),\Q_p/\Z_p\right).
\]
This describes the intended type of object; the local conditions defining the Selmer group in this problem are not supplied. Likewise, \(L_p(\phi_F,K[p^\infty])\) denotes the intended \(p\)-adic \(L\)-function, whose construction, normalization, and coefficient ring are not supplied. For an element \(f\) of a commutative ring, \((f)\) denotes its principal ideal. The characteristic ideal below is understood only after the relevant Iwasawa algebra, module, and hypotheses ensuring that it is defined have been fixed.

\textbf{Source-qualified target.} With all these objects interpreted as in the cited report, the extracted main-conjecture question is whether
\[
\bigl(L_p(\phi_F,K[p^\infty])\bigr)
=
\operatorname{char}_{\Lambda(J_K(\infty))}
 X(E/K[p^\infty]).
\]
The supplied correction reports that torsion of \(X(E/K[p^\infty])\) over the intended Iwasawa algebra and the ideal inclusion
\[
\bigl(L_p(\phi_F,K[p^\infty])\bigr)
\subseteq
\operatorname{char}_{\Lambda(J_K(\infty))}
 X(E/K[p^\infty])
\]
are already known under the source's hypotheses. Thus the remaining target in that corrected formulation is the reverse ideal inclusion
\[
\operatorname{char}_{\Lambda(J_K(\infty))}
 X(E/K[p^\infty])
\subseteq
\bigl(L_p(\phi_F,K[p^\infty])\bigr).
\]
These reports of known results are supplied source qualifications, not independently verified assertions here. The earlier formulation instead asks for torsion of \(X(E/k_\infty)\) over \(\Lambda(G)\) and equality of its characteristic ideal with a \(p\)-adic \(L\)-function ideal. The record does not identify \(k_\infty\) with \(K[p^\infty]\), or \(G\) with \(J_K(\infty)\). Consequently, the displayed target preserves the corrected extraction but is not a fully specified standalone conjecture until the missing definitions and hypotheses are recovered.
\subsection{Short English statement}
Under the listed hypotheses, the corrected source asks for equality between a \(p\)-adic \(L\)-function ideal and the characteristic ideal of a Selmer dual over a dihedral ring-class tower. It reports torsion and one ideal inclusion as known, leaving the reverse inclusion as the target. Several essential definitions and hypotheses are missing, so this is a source-qualified statement-only draft requiring mathematical review.
\subsection{Sources}
\begin{itemize}
\item Local UnsolvedMath registry, permanent ID \(30001276\), exact source problem code \texttt{catalogue entry}, title ``Iwasawa Main Conjecture for Dihedral Ring Class Towers,'' in the Oberwolfach Reports Open Problems collection. The supplied record gives proposed year 2009 and names no individual proposer.
\item Supplied report citation: \emph{Algebraische Zahlentheorie} (2009), Oberwolfach Reports, pp.~1667--1730, DOI \texttt{10.4171/owr/2009/30}. \url{https://doi.org/10.4171/owr/2009/30}. The collection background describes these reports as documenting workshops and open-problem discussions. The record attributes its correction to the report's dihedral main conjecture and Theorem~2.2; no relevant report excerpt or pinpoint page is supplied. This attribution has not been independently verified here.
\item Supplied literature source, labelled ``two-variable results toward the conjecture'': \url{https://arxiv.org/abs/1110.6197}.
\item Supplied literature source, labelled ``CM anticyclotomic Iwasawa theory'': \url{https://arxiv.org/abs/math/0302319}.
\item Supplied literature source, labelled ``original OWR report'': \url{https://doi.org/10.4171/owr/2009/30}.
\item Literature assessment dated exactly \texttt{2026-08-21}: the supplied curation labels the problem ``partially solved'' and states that CM cases and several one-variable specializations are known, while the full two-variable equality in the extracted setting is not established in general. This is supplied dated curation derived from the OWR collection, not an independently verified literature assessment. The record itself requests independent verification of time-sensitive status and cited claims. Its formulation label \texttt{corrected\_verified} is likewise supplied metadata, not verification performed for this draft.
\end{itemize}
\end{document}
